\documentclass[10pt,norsk,a4paper,usenames,dvipsnames]{article}
\usepackage[utf8]{inputenc}
\usepackage[T1]{fontenc}
\usepackage[norsk]{babel}
% - PDF-relatert
\usepackage{hyperref,pdfpages,hypcap,download}
\hypersetup{colorlinks=true,allcolors=.}
\newcommand\fhref[2]{%
	\href{#1}{#2}\footnote{\url{#1}}%
}
% Andre pakker
\usepackage[cm]{fullpage}
\usepackage{parskip,multicol,textcomp,amssymb,graphicx,color,enumitem,cleveref}
% - Korreksjon av fotnoter i seksjoner/overskrifter
\usepackage[stable]{footmisc}
% - Skrifttype
\usepackage[bitstream-charter]{mathdesign}
% - Kommentarer
\usepackage{comment}
\usepackage{epstopdf}
\usepackage[dvipsnames]{xcolor}
\usepackage{xfrac}
\usepackage{setspace}
\usepackage{etoolbox}



% Redefines the first level of enumerated lists.
\renewcommand{\theenumi}{\alph{enumi}}
\renewcommand{\labelenumi}{\theenumi)}

\epstopdfsetup{outdir=./}

\download[vedtekter.pdf]{protokoll/cyb-gf-protokoll-2021-vår.pdf}


\title{\huge Protokoll ordinær generalforsamling}
\author{\LARGE Cybernetisk Selskab}
\date{3. mai 2023}

\begin{document}


\maketitle



\begin{center}


\includegraphics[width=0.6\textwidth,height=0.6\textheight,keepaspectratio=true]{cyblogoa3.pdf}

\end{center}


\newpage


\tableofcontents

\section{Valg av møteleder}
Marius Torsvoll foreslås. Godkjent ved akklamasjon.

\section{Valg av referent}
Trine Høyås er foreslått. Godkjent ved akklamasjon.

\section{Valg av protokollunderskrivere}
Nikolas Hemer Martin og Martin Kristoffer Jørstad Heggem melder seg frivillig.

\section{Valg av tellekorps}
Nim og Eric Svebakk melder seg frivillig.

\section{Godkjenning av innkalling}
Godkjent ved akklamasjon.

\section{Godkjenning av dagsorden}
Godkjent ved akklamasjon.


\section{Semesterberetninger}
    \subsection{Semesterberetning ved leder}

    \begin{multicols}{2}


    Kjære medlemmer og interne,

    takk for ett flott semester så langt. Selv om semesteret fortsatt ikke er ferdig og vi går inn mot en travel eksamensperiode vil jeg gjerne dele noen av høydepunktene fra semesteret som har vært så langt.

    Siden semesterstartfesten for nesten 100 dager siden, har vi arrangert 53 større og mindre arrangementer, litt over ett arrangement hver andre dag og det er lagt opp til rundt 20 flere før semesteret er omme. 25 av disse arrangementene har vært utlån, 7 større samarbeidsarrangementer (som for eksempel ettermiddagen og WCFD), 16 ordinære torsdags- og fredagspubber, samt 4-5 internarrangementer. Dette er bare et utvalg av aktivitetene våre, da vi også har kosetirsdager og kafedrift. Vi er en travel forening hvor det alltid skjer noe.

    De interne har gjort en fantastisk jobb og stått på hardt med gjennomføringen av alle arrangementene våre, selv om det hender at gjestene uteblir og man kan slite med å få tak i nok funksjonærer. Dette er dessverre ikke bare en trend vi har sett, men også flere av de andre kjellerpubbene på Blindern. Både gjestene og de frivillige er vanskeligere å få tak i enn før. Vi kan ikke være sikre på hva som er årsakene, men flere ting kan trekkes fram som mulige forklaringer. For eksempel økende priser, lav studiestøtte og mer tidskrevende timeplaner. Til tross for disse utfordringene gjør de frivillige en flott innsats og stiller opp når det trengs.
    Det er heller ikke bare studentene som ser økende priser, men også vi som forening. For oss har for eksempel innkjøpsprisen på ett fat med pils økt med nesten 20\% siden januar. Dette har, som dere sikkert har sett, tvunget oss til å øke prisen på en pils fra 50/40 kr til 60/50 kr. Vi opplever også økte utgifter på grunn av prisøkninger på flere områder, men vi jobber for at vi skal ha ett lavt utgifstnivå slik at ikke studentene får en større støyt.

    Medlemstallet viser seg også å ha stabilisert seg etter korona og har lagt seg på normalen fra 2017-19. Vi har også gjort noen viktige fremskritt når det gjelder vedlikehold av lokalet vårt og oppdatering av rutiner, slik som brannrutinene og rutiner for opplæring av funksjonærer og mestere. Vi har blant annet fått ny steamer og prosjektor noe vi kan takke tilskudsordningene til velferdstinget for, samt administrasjonen på IFI.
    Vi har etablert god kontakt med flere viktige personer i administrasjonen, og vi jobber og har jobbet for å sikre at vi frivillige blir husket og mer inkludert. Dette blant annet ved å jobbe for at vi havner i rutinene og årshjulene deres da denne kontakten fort kan mistes og glemmes fra vår side ved styreendringer noe som gjør det særdeles viktig at begge parter følger opp. Vi har også prøvd å få en bedre kontakt med eiendomsavdelingen, som vi "leier" lokalet av, med varierende hell. Vi har for eksempel fortsatt ikke noen ny kjelleravtale etter at vi har purret siden august i fjor for å få på plass en ny avtale, men det er også flere andre barrierer vi har prøvd å bryte ned for å få gjøre det enklere for oss frivillige.

    Jeg vil avslutte med å takke alle medlemmer og frivillige for deres innsats og engasjement så langt i semesteret. Det er dere som gjør at foreningen går framover, og vi setter stor pris på deres engasjement og innsats. Til slutt vil jeg ønske dere alle lykke til med eksamensperioden og en velfortjent sommerferie når den tid kommer.

    \end{multicols}

    \textbf{Peter Hjelle Petersen-Øverleir}, \\
    Leder, \date{\emph{03. mai 2023}}

Ingen kommentarer eller spørsmål til beretningen.

    \subsection{Semesterberetning ved kjellermogul}
Hanne går gjennom sin semesterberetning. Ingen kommentarer eller spørsmål til beretningen.

\section{Økonomi}

    \subsection{Orientering}
    \subsection{Budsjett}
    \subsection{Regnskap}
    \subsection{Godkjenning av regnskap}

Økonomiansvarlig går gjennom økonomien og regnskapet.

Spørsmål til økonomiansvarlig:
Q: Hvor mye har vi hatt inn og ut sånn ca?
A: Usikker, ~1mill kanskje? Tar man med alt av støtte o.l. ~1,3m.

Forsamlingen godkjenner orienteringen.

\section{Kontingentfastsettelse}
    Hovedstyret foreslår å holde medlemskontingenten på kr.~50,-.

\textbf{Kontingenten forblir kr.~50,-.}

\section{Valg}

\begin{multicols}{2}

    \subsection{Hovedstyret}
        Man velges inn i hovedstyret for to semestre av gangen.
    Peter presenterer hvem og hva hovedstyret gjør.
        \subsubsection{Leder}
        Peter presenterer vervet.

        Ingen kandidater stiller i første omgang.

        Ingen kandidater stiller i andre eller tredje omgang.


        Peter stiller til slutt til vervet og presenterer seg selv.


        Q: Hva syns du om energiøl?
        A: Alkohol er downer, koffein er upper, de utligner seg selv.

        Q: Kommer du til å bli hele perioden?
        A: Jeg kan si allerede nå at jeg trekker meg til neste generalforsamling men blir ut neste semester.


        \textbf{Peter Hjelle Petersen-Øverleir velges til vervet ved akklamasjon}.

        \subsubsection{Nestleder}
        Peter presenterer vervet.

        Silje Helgesen stiller til valg og presenterer seg selv.

        Spørsmål fra salen:
        Q: Er du forberedt på å bli leder om ingen andre stiller?
        A: Nå er det en grunn til at jeg stiller til nestleder og ikke leder, men jeg skjønner jo at det står i vervbeskrivelsen at man er stedfortreder for leder.

        \textbf{Silje Helgesen velges til vervet ved akklamasjon}.

        % \subsubsection{Kjellermogul}


        \subsubsection{Økonomiansvarlig}
        Hanne presenterer vervet.

        Iver Nygård Juul stiller til valg, presenterer seg selv.

        Spørsmål fra salen:
        Q: Liker du pingviner?
        A: Ja

        Q: Kommer du til å fortsette i SPF?
        A: Ja

        Q: Kan vi få flere klistremerker?
        A: Ja

        Q: Har du andre foreningserfaringer?
        A: Er styremedlem i Mikro

        \textbf{Iver Nygård Juul velges til vervet ved akklamasjon.}

        % \subsubsection{Arrangementssjef}


        % \subsubsection{Internansvarlig}


        % \subsubsection{Promoteringssjef}


        \subsubsection{Rekrutteringsansvarlig}

        Peter presenterer vervet.

        To kandidater stiller til vervet og presenterer seg selv; Emma Fjeld Brovold og Knut Holtet

        Q: Er dere villige til å gå med pingvindrakt?
        A1: Ja
        A2: Jeg er villig til å svette.

        \textbf{Etter stemmmetelling er Emma Sofie Fjeld Brovold valgt til ny rekrutteringsansvarlig.}

        % \subsubsection{X-sjef}


    \subsection{Kjellerstyret} %TODO Oppdater listen med verv
        Alle verv som er til valg i kjellerstyret gjelder for ett semester av gangen. Med unntak av Kasserer som blir valgt inn for to semester.

Hanne presenterer kjellerstyret.


        \subsubsection{Kjellernestleder}
            % \textit{Vervet presenteres.}
        Hanne presenterer vervet.

        Knut Holtet stiller til vervet, og presenterer seg selv.

        Q: Hvis du har 100kr til å kjøpe inn snacks når alt er stengt og det er møte om 15 minutter?
        A: Funnet døgnåpen butikk. Hadde brukt alle pengene på smash.

        Q: Er du forberedt på å være møteleder om kjellermogul er syk?
        A: Ja

        Q: Relevant erfaring for vervet?
        A: Var veldig aktiv intern, sitter i styret i Defi. Vært med i styret i lokal fotballklubb.

        Q: Når begynte du på ifi?
        A: Nå i august.

        Q: Siden du er med i to foreninger, hvordan blir det med tiden fordelt mellom dem?
        A: Jeg har alltid svart at jeg ikke har noe liv så det er ikke noe problem.


        \textbf{Knut Holtet velges til vervet ved akklamasjon}.


        \subsubsection{Barsjef}
            % \textit{Vervet presenteres.}

        Mattis presenterer vervet.

        Ingen kandidater stiller i første omgang.
        Ingen kandidater stiller i andre omgang.



        \textbf{Julie Grihamar bestemmer seg for å stille, valgt til vervet ved akklamasjon. }


        \subsubsection{Kafésjef}
            % \textit{Vervet presenteres.}

        Martin presenterer vervet.

        Lise Chen stiller til vervet.

        Q: Var det du som dro i brannalarmen?
        A: Nei

        Q: Hvilken sorteringsalgoritme liker du for melk?
        A: FIFO så klart.

        Q: Favorittdrikk?
        A: Dobbel latte

        Q: Hvordan steamer man melk ordentlig?
        A Først så tar man steameren litt nedi så du får en brumming, og ikke sånn skriking[...]

        Q: Hva slags havremelk vil du kjøpe inn?
        A: Det vet jeg ikke, de er gode til hver sin ting.



        \textbf{Lise Chen velges til vervet ved akklamasjon.}

        % \subsubsection{Kasserer}
        %     % \textit{Vervet presenteres.}


        \subsubsection{Innkjøpsansvarlig}
            % \textit{Vervet presenteres.}
        Mick presenterer vervet og svarer på spørsmål.

        Q: Hvor ble det av Ginger Joe?
        A: Han ble borte.. noen liiker det, noen gjør det ikke.

        Q: Hvor ble det av toffee apple cideren?
        A: Det er mange grunner til at ting blir borte [...] Kollektivt liker folk mangocider.

        Q: Hva er favorittingen du har kjøpt inn?
        A: Mye nytt glutenfritt, mye mangocider. Ikke alt jeg har kjøpt inn er like bra men..

        Q: Hva anbefaler du neste innkjøpsansvarlig å kjøpe inn nå når det er tomt for mangocider?
        A: Det er ikke lenger tomt.



        Mick Guiotto Williams stiller til vervet.

        Spøsrmål fra salen:

        Q: Du sa du skulle kjøpe inn mindre, hva mener du med det?
        A: På ifi liker man å drikke tre ting og det er pils, mangocider og energiøl. Vi har hatt mye rart før, blir fort svinn. Vil ha mindre utvalg men raskere rullering på varer.

        Q: Kommer du til å sørge for at Cyb kommer på Vecturamesse?
        A: [Ja]

        Q: Kommer krigen mellom Cyb og RF om best utvalg til å legges død?
        A: [Ja]

        Q: Kommer Ginger Joe tilbake?
        A: Om ønskelig, vi har alkoholfri.


         \textbf{Mick Guiotto Williams velges til vervet ved akklamasjon.}


        \subsubsection{Teknisk sjef}
            % \textit{Vervet presenteres.}
        Elias presenterer vervet.

        Elias Hoel Birketvedt stiller til vervet og presenterer seg selv.

        Spørsmål fra salen:

        Q: Si at noen hypotetisk sett har lyst til å gå gjennom alt som henger i taket her i ferien, samt lage dokumentasjon til det er du interessert i det?
        A: Svaret er ja.

        Q: Har du mottatt flaskepost ennå?
        A: 1: aldri, 2: vi har kuttet flaskepost fra vervet da det aldri kommer flaskepost.

        Q: Plan for lys?
        A: Skal gå gjennom med Peter i sommeren og skal lage dokumentasjon på det.

        \textbf{Elias Hoel Birketvedt velges til vervet ved akklamasjon.}



        \subsubsection{Utlånsansvarlig}
            % \textit{Vervet presenteres.}
        Silje presenterer vervet.

        Q: Hva er timeslønnen?
        A: 250kr for SM, 200kr for funk.

        Q: Hva står SPF for?
        A: SPF står for studentkjellernes personalforening


        Julie og Ronja presenterer seg selv.

        Spørsmål fra salen:

        Q: Vil dere sagt at dere er flinke til å svare på mail og holde styr på datoer?
        A1: Ja
        A2: Ja, jeg har og funnet ut at google kalender er en god ting.

        Q: Er det noen av dere som vil bli leder i SPF?
        A1: Tror jeg måtte vært med på et møte først for å se hvordan det er, stiller ikke blind til ting
        A2: Virker som om Cyb allerede har gjort sin jobb der men potensielt

        Q: klarer dere å skjønne når dere har for mye å gjøre og evt si nei?
        A1: Ja
        A2: Ja, jeg vil si det. er det fullt så er det fullt.

        Q: Hvor lenge tenker dere å bli i Cyb?
        A1: Sånn som det går med IN1010 så..
        A2: jeg ser for meg å bare bli mer aktiv så vi får se.

        \textbf{Etter opptelling av stemmene har Ronja Karine Instanes blitt valgt til ny utlånsansvarlig.}

        \subsubsection{DJ-sjef}
            % \textit{Vervet presenteres.}
        Peter presenterer vervet.

        Mattis Stene-Johansen stiller til vervet og presenterer seg selv.

        Q: Kan du slutte å spille ABBA?
        A: Ja.. eller begrense det litt

        Q: Hvor viktig er transitions for deg?
        A: Veldig. Er lei av spotifylister

        Q: Favorittsang å spille?
        A: Har ingen favoritt, liker mye forskjellig musikk

        Q: Gleder du deg til å få nytt DJ-utstyr?
        A: Absolutt

        Q: [Har du tenkt til å gjøre noe med at det er min (Peter) sin jazzliste som er det vi spiller mest?]
        A: Vi må ha en slags musikkprofil, lage god stemning ikke bare jazz.

        Q: Hvordan ser du på rekruttering for DJ-gruppa?
        A: [Har lyst til å kapre alle technofans, alle som vil bidra for å lage en god dj-gruppe.]


\textbf{        Mattis Georg Nikolai Stene-Johansen velges til vervet ved akklamasjon.
}

        \subsubsection{Arrangementskoordinator}
            % \textit{Vervet presenteres.}

            Lina presenterer vervet.

             Miwa Sofie Meling Dimmen stiller til vervet og presenterer seg selv.


             Q: Hva var din favorittdrink?
             A: Jeg liker mange, men jeg tror en klassisk mojito er det jeg ville gått for.

             Q: Blir det torsdagspub i morgen?
             A: Vi snakka om det, vi vet ikke. Sliter med å finne SM og funker.

             Q: Har du god ide til temadrink?
             A: Jeg har tennkt på litt, men skal ikke si noe nå så det blir spennende senere.

\textbf{             Miwa Sofie Meling Dimmen velges til vervet ved akklamasjon.
}



\end{multicols}

\section{Vedtektsforslag}

     Peter presenterer de innsendte endringsforslagene. Ingen spørsmål fra salen.


    \subsection{§6 b - Endring}
    Endre:
    Kjellerstyret skal ha fra fem til ti personer, og ledes av kjellermogul med \textcolor{BrickRed}{barsjef} som stedfortreder.

    Til:
    Kjellerstyret skal ha fra fem til ti personer, og ledes av kjellermogul med \textcolor{ForestGreen}{kjellernestleder} som stedfortreder

    \textbf{Vedtektsendringen vedtas ved flertall.}


    \subsection{§6 c - Endring}
    Endre:
    Faste verv i kjellerstyret er kjellermogul, barsjef og kasserer. Øvrige verv defineres av kjellerstyret før generalforsamling.

    Til:
    Faste verv i kjellerstyret er kjellermogul, \textcolor{ForestGreen}{kjellernestleder}, barsjef og kasserer. Øvrige verv defineres av kjellerstyret før generalforsamling.

    \textbf{Vedtektsendringen vedtas ved flertall.}


    \subsection{§11 a - Presisering}
    Endre:
    Foreningens medlemmer jf. §2 kan fremme mistillitsforslag mot styremedlemmer og andre tillitsvalgte som er valgt i henhold til §8.

    Til:
    Foreningens medlemmer jf. §2 kan fremme mistillitsforslag mot styremedlemmer og andre tillitsvalgte som er valgt i henhold til §8\textcolor{ForestGreen}{, eller etterfylt etter §5 f) eller §6 e).}

    \textbf{Vedtektsendringen vedtas ved flertall.}

\section{Utdeling av utmerkelser}

Arkivar deler ut utmerkelser:

Pins:
\begin{itemize}
    \item Fredrik
    \item Silje
    \item Ingela
    \item Lina
    \item Mick
    \item Peter T.
    \item Philip
    \item Martin
\end{itemize}

Gullkort: \\
\textbf{Peter Hjelle Petersen-Øverleir
}
\section{Eventuelt}

Peter bestemmer seg for å stille til leder allikevel.

\newpage

\textbf{\huge{Hovedstyret består nå av:}}

\mbox{}


\begin{center}
\setstretch{2.0}
    \LARGE Leder -- Peter Hjelle Petersen-Øverleir
    \\ Nestleder -- Silje Helgesen
    \\ Økonomiansvarlig -- Iver Nygård Juul
    \\ Kjellermogul -- Hanne Alexandersdatter Hugem
    \\ Arrangementssjef -- Eric-Andreas Tepohe Urke Svebakk
    \\ Rekrutteringsansvarlig -- Emma Sofie Fjeld Brovold
    \\ Internansvarlig -- Kristian Tereshchenko Stillingen
    \\ Promoteringssjef -- Kavin Thas
    \\ X-sjef -- Kamil K\MakeLowercase{RZYŻANOWSKI}
\end{center}

\mbox{}
\\[24pt]

\begin{center}
\LARGE Protokollunderskrivere:\\[30pt]

\begin{multicols}{2}
    \Large
    Martin Kristoffer Jørstad Heggem\\[48pt]
    \rule{220pt}{1pt}

    Nikolas Hemer Martin\\[48pt]
    \rule{220pt}{1pt}


\end{multicols}
\end{center}


% \addcontentsline{toc}{part}{Vedlegg}

% \newpage
% \phantomsection{}
% \addcontentsline{toc}{section}{Vedtekter for Cybernetisk Selskab} % chktex
% \includepdf[pages=-]{cyb_vedtekter.pdf}

\end{document}
